% Included by main.tex; the preamble and references are shared.
\section{Boundary kernels and the multiplier identity}\label{sec:boundary-kernels}

\subsection{A subcritical disk potential}

Let $\Omega\subset\mathbb R^2$ be a nonempty bounded simply connected
open set with smooth boundary, and fix the smooth conformal
diffeomorphism $\Phi:\overline{\mathbb D}\to\overline\Omega$ constructed
in \cref{lem:conformal}. Here
$\mathbb D=\{z\in\mathbb C:|z|<1\}$ and $S^1=\partial\mathbb D$.
Write $\lambda_D$ and $\lambda_N$ for the first Dirichlet and first
positive Neumann eigenvalues of $\Omega$, respectively. Define
\begin{equation}\label{eq:conformal-measure}
 dq(x)=\lambda_N|\Phi'(x)|^2\,dx.
\end{equation}
Its density is bounded and positive. We use real scalar functions and
write $\langle f,g\rangle_q=\int_{\mathbb D}fg\,dq$ and
$\|f\|_q^2=\langle f,f\rangle_q$.

For $\eta\in H_0^1(\mathbb D)$, put $u=\eta\circ\Phi^{-1}$.
The smooth diffeomorphism preserves zero trace, so
$u\in H_0^1(\Omega)$. Conformality gives
$D\Phi=|\Phi'|R$, where $R$ is a rotation, and
$\det D\Phi=|\Phi'|^2$. Thus the chain rule and change of variables give
\[
 \int_\Omega|\nabla u|^2\,dy
 =\int_{\mathbb D}|\nabla\eta|^2\,dx.
\]
These identities hold first for smooth functions and then by density
in $H_0^1$. The Dirichlet inequality on $\Omega$ therefore gives
\[
 \|\eta\|_q^2
 =\lambda_N\int_\Omega u^2\,dy
 \le\frac{\lambda_N}{\lambda_D}\int_\Omega|\nabla u|^2\,dy
 =\frac{\lambda_N}{\lambda_D}
       \int_{\mathbb D}|\nabla\eta|^2\,dx.
\]
The spectral comparison \cref{lem:spectral-gap} gives
$\lambda_N<\lambda_D$. Hence
\begin{equation}\label{eq:weighted-dirichlet}
 \|\eta\|_q^2\le d\int_{\mathbb D}|\nabla\eta|^2\,dx
 \qquad(\eta\in H_0^1(\mathbb D)),
 \qquad d=\frac{\lambda_N}{\lambda_D}<1.
\end{equation}

For the scalar kernel construction, we allow any measure $dq=a\,dx$
with bounded nonnegative density satisfying
\cref{eq:weighted-dirichlet} for some fixed $0<d<1$.
The conformal measure \cref{eq:conformal-measure} is one such measure;
the zero measure is also allowed.

The Dirichlet Green function of the disk is
\begin{equation}\label{eq:disk-green}
 G(x,y)=\frac1{2\pi}\log\left|
                  \frac{1-\overline{x}y}{x-y}\right|
 =\frac1{4\pi}\log\left(
       1+\frac{(1-|x|^2)(1-|y|^2)}{|x-y|^2}\right),
 \qquad x\ne y.
\end{equation}
It is nonnegative and symmetric, vanishes when either variable lies
on $S^1$, and has logarithmic singularity
$-(2\pi)^{-1}\log|x-y|$. Its diagonal value is set to $+\infty$;
the diagonal has measure zero in all integrals below.
For fixed $x\in\mathbb D$, it is harmonic in $y$ away from $x$:
both logarithms in the first formula are logarithms of moduli of
nonvanishing holomorphic functions there.

\begin{lemma}\label[lemma]{lem:green-operator}
The operator
\[
 Tf(x)=\int_{\mathbb D}G(x,y)f(y)\,dq(y)
\]
is bounded and self-adjoint on $L^2(q)$, with $\|T\|\le d<1$.
For each $f\in L^2(q)$, the displayed integral is the continuous
interior representative of the weak Dirichlet solution of
$-\Delta\psi=af$.
For every measure $0\le\rho\le q$, the operator
\[
 T_\rho f(x)=\int_{\mathbb D}G(x,y)f(y)\,d\rho(y)
\]
has the same properties on $L^2(\rho)$, with $\|T_\rho\|\le d$.
\end{lemma}
\begin{proof}
We first justify the Green representation for a source
$b\in L^2(\mathbb D;dx)$. The Poincar\'e inequality and Lax--Milgram
\cite[Theorem 4.20, p.~103]{Hunter}
give a unique $u\in H_0^1(\mathbb D)$ satisfying
\begin{equation}\label{eq:weak-poisson}
 \int_{\mathbb D}\nabla u\cdot\nabla\eta\,dx
 =\int_{\mathbb D}b\eta\,dx
 \qquad(\eta\in H_0^1(\mathbb D)).
\end{equation}
If $b\in C_c^\infty(\mathbb D)$, this solution is smooth up to the
boundary. Apply Green's second identity to $u$ and $G(x,\cdot)$ on
$\mathbb D\setminus\overline{B_\varepsilon(x)}$.
The outer boundary term is zero. On the inner circle, whose outward
normal points towards $x$,
\[
 G(x,y)=-\frac1{2\pi}\log\varepsilon+O(1),
 \qquad \partial_\nu G(x,y)=\frac1{2\pi\varepsilon}+O(1).
\]
Consequently the integral of
$G\partial_\nu u-u\partial_\nu G$ tends to $-u(x)$.
Since $\Delta_yG(x,y)=0$ away from $x$ and $\Delta u=-b$,
letting $\varepsilon\downarrow0$ gives
\begin{equation}\label{eq:green-representation}
 u(x)=\int_{\mathbb D}G(x,y)b(y)\,dy.
\end{equation}

For general $b\in L^2(dx)$, choose $b_n\in C_c^\infty(\mathbb D)$
converging to $b$ in $L^2$, and let $u_n$ be their weak solutions.
Testing the difference equation and using Poincar\'e gives
\[
 \|\nabla(u_n-u)\|_2^2
 \le C\|b_n-b\|_2\|\nabla(u_n-u)\|_2.
\]
Thus $u_n\to u$ in $H_0^1$, and a subsequence converges almost
everywhere. For each fixed $x\in\mathbb D$, the logarithmic pole is
square integrable, so $G(x,\cdot)\in L^2(dx)$ and
\[
 \left|\int_{\mathbb D}G(x,y)(b_n-b)(y)\,dy\right|
 \le\|G(x,\cdot)\|_2\|b_n-b\|_2\longrightarrow0.
\]
Passing to the limit in \cref{eq:green-representation} proves that
its right-hand side equals $u$ almost everywhere.

To prove continuity of this representative, fix $K\Subset\mathbb D$
and $r<1$ with $|x|\le r$ on $K$. For $x,x_0\in K$,
\[
 \left|\log|1-\overline{x}y|
           -\log|1-\overline{x_0}y|\right|
 \le\frac{|x-x_0|}{1-r}\qquad(y\in\mathbb D).
\]
For the singular part choose $\chi\in C_c^\infty(\mathbb R^2)$
equal to one on $\{|z|\le2\}$ and set $F(z)=\chi(z)\log|z|$.
Then $F\in L^2(\mathbb R^2)$, since
$\int_0^1 t|\log t|^2\,dt<\infty$, and
\[
 \|\log|x-\cdot|-\log|x_0-\cdot|\|_{L^2(\mathbb D)}
 \le\|F(x-\cdot)-F(x_0-\cdot)\|_{L^2(\mathbb R^2)}
 \longrightarrow0.
\]
Here continuity of translations in $L^2$ follows by approximation
by smooth compactly supported functions: translation preserves the
$L^2$ norm, and the translated approximants converge in $L^2$.
Together with \cref{eq:disk-green}, these estimates give
$\|G(x,\cdot)-G(x_0,\cdot)\|_2\to0$. Hence
\[
 \left|\int_{\mathbb D}(G(x,y)-G(x_0,y))b(y)\,dy\right|
 \le\|G(x,\cdot)-G(x_0,\cdot)\|_2\|b\|_2\longrightarrow0.
\]
This proves the asserted interior continuity.

Now let $f\in L^2(q)$. The source $af$ belongs to $L^2(dx)$ because
\[
 \int_{\mathbb D}|af|^2\,dx
 \le\|a\|_\infty\int_{\mathbb D}|f|^2\,dq.
\]
Changing $f$ on a $q$-null set does not change $af$ almost everywhere
with respect to $dx$, so the integral defines a linear map on $L^2(q)$.
The representation just proved identifies $\psi=Tf$ with the weak
solution for $b=af$. Taking $\eta=\psi$ in \cref{eq:weak-poisson}
gives
\[
 \int_{\mathbb D}|\nabla\psi|^2\,dx
 =\int_{\mathbb D}af\psi\,dx=\langle\psi,f\rangle_q.
\]
Thus \cref{eq:weighted-dirichlet} and Cauchy--Schwarz imply
\[
 \|\psi\|_q^2
 \le d\int_{\mathbb D}|\nabla\psi|^2\,dx
 =d\langle\psi,f\rangle_q
 \le d\|\psi\|_q\|f\|_q.
\]
Dividing when $\|\psi\|_q>0$, with the zero case immediate, proves
$\|T\|\le d$.

For self-adjointness, \cref{eq:disk-green} gives
$G(x,y)=G(y,x)$ and
\[
 0\le G(x,y)\le\frac1{2\pi}\log\frac2{|x-y|}.
\]
The square of the upper bound is integrable on
$\mathbb D\times\mathbb D$; boundedness of $a$ therefore gives
$G\in L^2(q\times q)$. For $f,g\in L^2(q)$, Cauchy--Schwarz on the
product space gives
\[
 \iint G(x,y)|f(y)g(x)|\,dq(y)\,dq(x)
 \le\|G\|_{L^2(q\times q)}\|f\|_q\|g\|_q<\infty.
\]
Fubini and symmetry now yield
\[
 \langle Tf,g\rangle_q
 =\iint G(x,y)f(y)g(x)\,dq(y)\,dq(x)
 =\langle f,Tg\rangle_q.
\]
Since $T$ is bounded and defined on all of $L^2(q)$, it is self-adjoint.

Finally, $0\le\rho\le q$ means $\rho(E)\le q(E)$ for every
measurable set $E$. It follows that $d\rho=c\,dx$ with
$0\le c\le a$ almost everywhere, so $c$ is bounded. Moreover,
\[
 \|\eta\|_{L^2(\rho)}^2
 \le\|\eta\|_q^2
 \le d\int_{\mathbb D}|\nabla\eta|^2\,dx
 \qquad(\eta\in H_0^1(\mathbb D)).
\]
This is the estimate needed to bound the solution operator on
$L^2(\rho)$. Indeed, for $f\in L^2(\rho)$ the source $cf$ lies in
$L^2(dx)$, and its weak solution $\psi=T_\rho f$ satisfies
\[
 \|\psi\|_{L^2(\rho)}^2
 \le d\int_{\mathbb D}|\nabla\psi|^2\,dx
 =d\int_{\mathbb D}f\psi\,d\rho
 \le d\|f\|_{L^2(\rho)}\|\psi\|_{L^2(\rho)}.
\]
Thus $\|T_\rho\|\le d$. The same Green representation proves
interior continuity, and the same Fubini argument with $c$ in place
of $a$ proves self-adjointness.
\end{proof}

For Lipschitz boundary data $h:S^1\to\mathbb R$, define the Poisson
extension
\begin{equation}\label{eq:poisson-extension}
 Ah(x)=\int_{S^1}P_s(x)h(s)\,ds,
 \qquad P_s(x)=\frac{1-|x|^2}{2\pi|s-x|^2},
\end{equation}
where $ds$ is arclength on the unit circle. To justify its Sobolev
regularity, write
\[
 h(e^{i\theta})=\sum_{n\in\mathbb Z}h_ne^{in\theta},
 \qquad
 Ah(re^{i\theta})=\sum_{n\in\mathbb Z}h_nr^{|n|}e^{in\theta}.
\]
Lipschitz regularity gives $h\in H^1(S^1)$, so
$\sum_n(1+n^2)|h_n|^2<\infty$. Parseval's identity and integration
in polar coordinates give
\[
 \|Ah\|_{L^2(\mathbb D)}^2
 =\pi\sum_{n\in\mathbb Z}\frac{|h_n|^2}{|n|+1},
 \qquad
 \|\nabla Ah\|_{L^2(\mathbb D)}^2
 =2\pi\sum_{n\in\mathbb Z}|n|\,|h_n|^2.
\]
These identities hold first for finite Fourier sums; their harmonic
extensions converge in $H^1(\mathbb D)$ by the displayed bounds.
Their boundary values converge to $h$ in $H^1(S^1)$, so continuity
of the trace shows that $\operatorname{Tr}Ah=h$. Harmonicity also
passes to the limit, giving
\[
 \int_{\mathbb D}\nabla Ah\cdot\nabla\eta\,dx=0
 \qquad(\eta\in H_0^1(\mathbb D)).
\]

We seek $W_h\in H^1(\mathbb D)$ with trace $h$ satisfying
\begin{equation}\label{eq:weak-potential-solution}
 \int_{\mathbb D}\nabla W_h\cdot\nabla\eta\,dx
 =\int_{\mathbb D}aW_h\eta\,dx
 \qquad(\eta\in H_0^1(\mathbb D)).
\end{equation}
This is the weak meaning of $-\Delta W_h\,dx=W_h\,dq$.
Existence and uniqueness follow from Lax--Milgram
\cite[Theorem 4.20, p.~103]{Hunter} after writing
$W_h=Ah+c$, $c\in H_0^1(\mathbb D)$. Indeed, the continuous form
\[
 B(u,\eta)=\int_{\mathbb D}\nabla u\cdot\nabla\eta\,dx
           -\int_{\mathbb D}au\eta\,dx
\]
satisfies
$B(c,c)\ge(1-d)\|\nabla c\|_2^2$ by
\cref{eq:weighted-dirichlet}, and the correction equation is
$B(c,\eta)=\int_{\mathbb D}aAh\,\eta\,dx$.

The solution has the explicit representation
\begin{equation}\label{eq:potential-extension}
 W_h=Ah+T(I-T)^{-1}Ah.
\end{equation}
To verify it, boundedness of $a$ gives $Ah\in L^2(q)$, and
\cref{lem:green-operator} gives
$(I-T)^{-1}=\sum_{n\ge0}T^n$ in operator norm on $L^2(q)$.
Put $v=(I-T)^{-1}Ah$ and $c=Tv$. Then
$c\in H_0^1(\mathbb D)$, and the weak Poisson equation gives
\[
 \int_{\mathbb D}\nabla c\cdot\nabla\eta\,dx
 =\int_{\mathbb D}av\eta\,dx
 \qquad(\eta\in H_0^1(\mathbb D)).
\]
The identity $(I-T)v=Ah$ implies $v=Ah+c=W_h$ in $L^2(q)$.
Consequently harmonicity of $Ah$ yields
\[
 \int_{\mathbb D}\nabla W_h\cdot\nabla\eta\,dx
 =\int_{\mathbb D}\nabla c\cdot\nabla\eta\,dx
 =\int_{\mathbb D}aW_h\eta\,dx.
\]
Also $\operatorname{Tr}W_h=h$, since $c$ has zero trace.
For a direct uniqueness check, the difference $z$ of two solutions
with this trace lies in $H_0^1(\mathbb D)$. Testing its equation by
$z$ and using \cref{eq:weighted-dirichlet} gives
\[
 \|\nabla z\|_2^2=\|z\|_q^2\le d\|\nabla z\|_2^2.
\]
Since $d<1$, $\nabla z=0$, and Poincar\'e gives $z=0$.
For vector boundary data all constructions and equations are
understood componentwise.

\subsection{Positive kernels and their estimates}

For $s\in S^1$, define $U_{0,s}=P_s$ and recursively set
\[
 U_{n+1,s}(x)=\int_{\mathbb D}G(x,y)U_{n,s}(y)\,dq(y).
\]
The nonnegative integrands define measurable functions with values in
$[0,\infty]$. We write $T^nP_s=U_{n,s}$; this definition does not
require $P_s\in L^2(q)$. Finiteness is established below.
For $p\in\mathbb D$ and distinct $s,t\in S^1$, define
\begin{align}
 K(p,s)&=P_s(p)+\sum_{n\ge1}(T^nP_s)(p),
                    \label{eq:solution-kernel}\\
 R(s,t)&=\sum_{n\ge0}\int_{\mathbb D}P_sT^nP_t\,dq,
                    \label{eq:interaction-kernel}\\
 N(s,t)&=N_0(s,t)+R(s,t),
 \qquad N_0(s,t)=\frac1{\pi|s-t|^2}.
                    \label{eq:boundary-kernel}
\end{align}
For integrations we also define $R$ on the diagonal by the same
nonnegative series, allowing the value $+\infty$.

\begin{lemma}\label[lemma]{lem:kernel-estimates}
The kernels in \cref{eq:solution-kernel,eq:interaction-kernel,eq:boundary-kernel}
are finite at the stated points. For fixed $p\in\mathbb D$,
$K(p,\cdot)$ is positive and bounded. The kernel $R$ is symmetric,
nonnegative, and integrable on $S^1\times S^1$, with
\begin{equation}\label{eq:interaction-mass}
 \iint_{S^1\times S^1}R(s,t)\,ds\,dt
 \le\frac{q(\mathbb D)}{1-d}.
\end{equation}
For Lipschitz scalar or vector boundary data $h,k$,
\begin{align}
 W_h(p)&=\int_{S^1}K(p,s)h(s)\,ds,
                         \label{eq:solution-representation}\\
 \langle Ah,(I-T)^{-1}Ak\rangle_q
 &=\iint_{S^1\times S^1}R(s,t)h(s)\cdot k(t)\,ds\,dt.
                         \label{eq:interaction-representation}
\end{align}
For scalar data the dot denotes multiplication; for vector data the
inner product on the left is the sum of the component inner products.
\end{lemma}
\begin{proof}
Constants below are independent of interior and boundary points,
unless stated otherwise. The explicit formulas give
\begin{equation}\label{eq:elementary-kernel-bounds}
 P_s(y)\le\frac1{\pi|y-s|},
 \qquad 0\le G(x,y)\le\frac1{2\pi}\log\frac2{|x-y|}.
\end{equation}
For the first bound use $1-|y|^2\le2|y-s|$; for the second use
$|1-\overline{x}y|\le2$. Translating the centres and integrating
over a disk of radius two gives
\[
 \sup_{s\in S^1}\|P_s\|_{L^{3/2}(dx)}<\infty,
 \qquad
 \sup_{x\in\mathbb D}\|G(x,\cdot)\|_{L^3(dx)}<\infty,
\]
since $\int_0^2 r^{-3/2}r\,dr$ and
$\int_0^2(\log(2/r))^3r\,dr$ are finite.

Define the Green potential $F_s=TP_s$. We use it as the starting
point for the $L^2(q)$ contraction estimate. Writing $dq=a\,dy$,
H\"older's inequality with exponents $3$ and $3/2$ gives
\begin{equation}\label{eq:first-green-bound}
 F_s(x)=\int_{\mathbb D}G(x,y)P_s(y)a(y)\,dy
 \le\|a\|_\infty\|G(x,\cdot)\|_{L^3(dx)}
                   \|P_s\|_{L^{3/2}(dx)}\le C.
\end{equation}
Thus Green integration regularises the singularity of $P_s$ at
$s$, producing a bounded function, with
$\|F_s\|_q\le Cq(\mathbb D)^{1/2}$. The recursively defined
integrals satisfy $T^nP_s=T^{n-1}F_s$ for $n\ge1$, with the
right-hand side interpreted as the $L^2(q)$ operator iterates and,
for $n\ge2$, their Green representatives.
Since $q$ has bounded density, the uniform $L^3(dx)$ bound above
also gives $\sup_p\|G(p,\cdot)\|_q<\infty$. Hence, for $n\ge2$,
\begin{equation}\label{eq:geometric-kernel-bound}
 (T^nP_s)(p)
 =\langle G(p,\cdot),T^{n-2}F_s\rangle_q
 \le\|G(p,\cdot)\|_q d^{n-2}\|F_s\|_q
 \le C d^{n-2}.
\end{equation}
Together with \cref{eq:first-green-bound}, this proves convergence
of the series for $K$. For fixed $p$, the function $P_s(p)$ is
positive and bounded as $s$ varies, so the same holds for $K(p,s)$.

For the zeroth term of $R$, fix $s\ne t$ and put
$\delta=|s-t|/3$. On $\mathbb D\cap B_\delta(s)$,
$|y-t|\ge2\delta$, and \cref{eq:elementary-kernel-bounds} gives
\[
 P_s(y)P_t(y)\le\frac1{2\pi^2\delta|y-s|}.
\]
This is integrable because
$\int_{B_\delta(s)}|y-s|^{-1}\,dy=2\pi\delta$.
The same argument applies near $t$. Outside these two disks,
$P_sP_t\le(\pi^2\delta^2)^{-1}$. Thus $P_sP_t\in L^1(dx)$,
and boundedness of $a$ gives $\int P_sP_t\,dq<\infty$.
Also \cref{eq:elementary-kernel-bounds} gives
$\sup_s\int_{\mathbb D}P_s\,dq\le C$. Therefore the explicit
bounds \cref{eq:first-green-bound,eq:geometric-kernel-bound} imply
\[
 \int_{\mathbb D}P_sTP_t\,dq\le C,
 \qquad
 \int_{\mathbb D}P_sT^nP_t\,dq\le C d^{n-2}\quad(n\ge2).
\]
This proves finiteness of $R(s,t)$, and hence of $N(s,t)$, for
distinct $s,t$.

For $n\ge1$, the $n$th term of $R$ is the nonnegative chain integral
\[
 \int_{\mathbb D^{n+1}}
 P_s(y_0)\left(\prod_{j=0}^{n-1}G(y_j,y_{j+1})\right)
 P_t(y_n)\,dq(y_0)\cdots dq(y_n).
\]
Reversing the variables and using symmetry of $G$ proves symmetry
in $s,t$; the zeroth term is also symmetric.
The Poisson kernel satisfies $\int_{S^1}P_s(y)\,ds=1$.
Tonelli's theorem, applied to the nonnegative chains and series,
therefore gives
\[
 \iint R(s,t)\,ds\,dt
 =\sum_{n\ge0}\langle1,T^n1\rangle_q
 \le\sum_{n\ge0}d^n\|1\|_q^2
 =\frac{q(\mathbb D)}{1-d}.
\]
This proves \cref{eq:interaction-mass}. Values on the diagonal
$s=t$ do not affect these integrals, since it has zero product
arclength measure.

We next prove \cref{eq:solution-representation}. The Neumann series
in \cref{eq:potential-extension} gives
\[
 W_h(p)=Ah(p)+\sum_{n\ge1}T^nAh(p).
\]
Indeed, the series for $(I-T)^{-1}Ah$ converges in $L^2(q)$,
and evaluation after applying $T$ is bounded by
$|Tf(p)|\le\|G(p,\cdot)\|_q\|f\|_q$.
For nonnegative scalar $h$, inserting \cref{eq:poisson-extension}
into the iterated integrals and applying Tonelli gives
\[
 T^nAh(p)=\int_{S^1}(T^nP_s)(p)h(s)\,ds
 \qquad(n\ge0).
\]
Summing gives \cref{eq:solution-representation}.
For signed bounded $h$, all exchanges are absolutely justified by
\[
 \int_{S^1}K(p,s)|h(s)|\,ds
 \le\|h\|_\infty\int_{S^1}K(p,s)\,ds<\infty,
\]
using boundedness of $K(p,\cdot)$.

Finally, the operator Neumann series gives
\[
 \langle Ah,(I-T)^{-1}Ak\rangle_q
 =\sum_{n\ge0}\langle Ah,T^nAk\rangle_q.
\]
For nonnegative scalar $h,k$, their Poisson representations and
Tonelli yield
\[
 \langle Ah,T^nAk\rangle_q
 =\iint h(s)k(t)
      \left(\int_{\mathbb D}P_sT^nP_t\,dq\right)\,ds\,dt.
\]
Summing proves \cref{eq:interaction-representation}. For signed
bounded data the required absolute domination follows from
\cref{eq:interaction-mass}:
\[
 \iint R(s,t)|h(s)k(t)|\,ds\,dt
 \le\|h\|_\infty\|k\|_\infty\frac{q(\mathbb D)}{1-d}<\infty.
\]
The vector formulas follow by applying the scalar formulas to each
component and summing.
\end{proof}

\subsection{The boundary energy}

We now return to the conformal measure \cref{eq:conformal-measure}.
A vector boundary function $h:S^1\to\mathbb R^2$ is called tangent
if $h(s)$ is tangent to $\partial\Omega$ at $\Phi(s)$.
Write $h=(h_1,h_2)$ in fixed Cartesian coordinates, and define
$W_h=(W_{h_1},W_{h_2})$, where each component is the scalar solution
in \cref{eq:potential-extension}. Set
\[
 X_h(y)=W_h(\Phi^{-1}(y)),\qquad y\in\Omega.
\]
Composition with the smooth diffeomorphism gives
$X_h\in H^1(\Omega;\mathbb R^2)$, and its trace satisfies
$\operatorname{Tr}X_h(\Phi(s))=h(s)$. Thus tangent data give tangent
trace on $\partial\Omega$.

As in the preceding section on tangent vector fields, write
\[
 Q(X)=\int_\Omega\bigl(|\operatorname{div}X|^2
                         +|\operatorname{curl}X|^2\bigr)\,dy,
 \qquad \operatorname{curl}X=\partial_1X_2-\partial_2X_1,
\]
and define
\[
 E(h)=Q(X_h)-\lambda_N\|X_h\|_{L^2(\Omega)}^2,
 \qquad E(h,k)=\frac14\bigl(E(h+k)-E(h-k)\bigr).
\]
The vector variational principle \cref{prop:vector-variational}
gives $E(h)\ge0$ for tangent
Lipschitz data, with equality exactly when $X_h=\nabla v$ for some
$v$ in the real Neumann eigenspace $V$ for $\lambda_N$.
Let $\nu$ be the outward unit normal to $\partial\Omega$,
$\tau=(-\nu_2,\nu_1)$, and
$\kappa=\det(\tau,\partial_\tau\tau)$, as in that section.

\begin{lemma}\label[lemma]{lem:boundary-energy}
For tangent Lipschitz data $h,k:S^1\to\mathbb R^2$,
\begin{equation}\label{eq:boundary-energy}
 \begin{aligned}
 E(h,k)
 &=\frac12\iint_{S^1\times S^1}
       N_0(s,t)(h(s)-h(t))\cdot(k(s)-k(t))\,ds\,dt\\
 &\quad-\iint_{S^1\times S^1}R(s,t)h(s)\cdot k(t)\,ds\,dt
       +\int_{S^1}|\Phi'|(\kappa\circ\Phi)h\cdot k\,ds.
 \end{aligned}
\end{equation}
\end{lemma}
\begin{proof}
Polarising the curvature identity \cref{lem:curvature-identity}
for tangent vector fields gives
\[
 E(h,k)=\int_\Omega\nabla X_h:\nabla X_k\,dy
       -\lambda_N\int_\Omega X_h\cdot X_k\,dy
       +\int_{\partial\Omega}\kappa\,
                 \operatorname{Tr}X_h\cdot\operatorname{Tr}X_k\,d\sigma,
\]
where $\nabla X_h:\nabla X_k
=\sum_{j=1}^2\nabla(X_h)_j\cdot\nabla(X_k)_j$.
For each Cartesian component, conformality and change of variables
give
\[
 \int_\Omega\nabla(X_h)_j\cdot\nabla(X_k)_j\,dy
 =\int_{\mathbb D}\nabla W_{h_j}\cdot\nabla W_{k_j}\,dx.
\]
Indeed, the inverse derivative contributes $|\Phi'|^{-2}$ to the
gradient product, the rotation preserves dot products, and the
area Jacobian is $|\Phi'|^2$. Also
\[
 \lambda_N\int_\Omega X_h\cdot X_k\,dy
 =\int_{\mathbb D}W_h\cdot W_k\,\lambda_N|\Phi'|^2\,dx
 =\langle W_h,W_k\rangle_q.
\]
For the boundary integral, the parametrisation
$\theta\mapsto\Phi(e^{i\theta})$ has speed
\[
 |\Phi'(e^{i\theta})ie^{i\theta}|=|\Phi'(e^{i\theta})|.
\]
Thus $d\sigma=|\Phi'(s)|\,ds$, and the trace identities give
\[
 \int_{\partial\Omega}\kappa\,
           \operatorname{Tr}X_h\cdot\operatorname{Tr}X_k\,d\sigma
 =\int_{S^1}|\Phi'|(\kappa\circ\Phi)h\cdot k\,ds.
\]

It remains to express the two integrals over $\mathbb D$ in terms
of boundary data. The difference $c_h=W_h-Ah$ has zero trace.
Testing the equation for each component of $W_k$ by the corresponding
component of $c_h$ gives
\[
 \int_{\mathbb D}\nabla c_h:\nabla W_k\,dx
 =\langle c_h,W_k\rangle_q.
\]
Consequently
\[
 \int_{\mathbb D}\nabla W_h:\nabla W_k\,dx
       -\langle W_h,W_k\rangle_q
 =\int_{\mathbb D}\nabla Ah:\nabla W_k\,dx
       -\langle Ah,W_k\rangle_q.
\]
Harmonicity of $Ah$ and zero trace of $W_k-Ak$ give
$\int_{\mathbb D}\nabla Ah:\nabla(W_k-Ak)\,dx=0$.
Since $W_k=(I-T)^{-1}Ak$ in $L^2(q)$, the last expression is
\[
 \int_{\mathbb D}\nabla Ah:\nabla Ak\,dx
       -\langle Ah,(I-T)^{-1}Ak\rangle_q.
\]
The second term is the integral with $R$ in
\cref{eq:interaction-representation}.

For the harmonic term we prove
\begin{equation}\label{eq:harmonic-boundary-energy}
 \int_{\mathbb D}\nabla Ah:\nabla Ak\,dx
 =\frac1{2\pi}\iint_{S^1\times S^1}
       \frac{(h(s)-h(t))\cdot(k(s)-k(t))}{|s-t|^2}\,ds\,dt.
\end{equation}
First consider scalar Fourier modes $e_m(s)=s^m$, $m\in\mathbb Z$.
For this calculation only, use complex functions and conjugate the
second factor in each energy pairing. The harmonic extension is
$Ae_m(re^{i\theta})=r^{|m|}e^{im\theta}$, and direct integration
gives
\[
 \int_{\mathbb D}|\nabla Ae_m|^2\,dx=2\pi|m|.
\]
For $m\ge1$,
\[
 \frac{s^m-t^m}{s-t}=\sum_{j=0}^{m-1}s^{m-1-j}t^j.
\]
The summands are orthogonal in $L^2(S^1\times S^1)$ and each has
squared norm $(2\pi)^2$. Hence
\[
 \iint\frac{|s^m-t^m|^2}{|s-t|^2}\,ds\,dt=(2\pi)^2m.
\]
Complex conjugation gives the same formula with $|m|$ for negative
$m$, and both energies vanish for $m=0$. After multiplication by
$1/(2\pi)$, this agrees with the disk energy.

Different modes have zero cross terms in the disk integral by
angular orthogonality. For the boundary integral, set
\[
 J_{m\ell}=\iint
 \frac{(s^m-t^m)\overline{(s^\ell-t^\ell)}}{|s-t|^2}\,ds\,dt.
\]
Rotating both variables by $e^{i\alpha}$ preserves arclength and
the denominator, so
$J_{m\ell}=e^{i(m-\ell)\alpha}J_{m\ell}$ for every $\alpha$.
Thus $J_{m\ell}=0$ when $m\ne\ell$.
The matching mode energies and vanishing cross terms prove
\cref{eq:harmonic-boundary-energy} for real scalar trigonometric
polynomials, and then for smooth data by convergence of their
Fourier sums in $C^1(S^1)$.

For Lipschitz data, smooth periodic convolution gives approximants
converging in $H^1(S^1)$ with uniformly bounded Lipschitz constants.
The Poisson extensions converge in $H^1(\mathbb D)$ by the Fourier
energy formulas above. On the boundary, the integrands are dominated
by a constant, since
\[
 \frac{|h(s)-h(t)|\,|k(s)-k(t)|}{|s-t|^2}
 \le\operatorname{Lip}(h)\operatorname{Lip}(k).
\]
Dominated convergence therefore proves the scalar identity for
Lipschitz data. Summing over the two components proves
\cref{eq:harmonic-boundary-energy} for vectors. Since
$N_0(s,t)=1/(\pi|s-t|^2)$, combining this identity with the preceding
calculations gives \cref{eq:boundary-energy}.
\end{proof}

Fix a nonzero $u\in V$, and define
\begin{equation}\label{eq:gradient-boundary-data}
 g(s)=(\nabla u)(\Phi(s)),\qquad s\in S^1.
\end{equation}
Elliptic regularity
\cite[Chapter III, Theorems 6.4 and 6.5]{Showalter}
gives $u\in C^\infty(\overline\Omega)$, so
$g$ is smooth. It is tangent by the Neumann condition
$\partial_\nu u=0$. For each Cartesian derivative
$v_j=\partial_j u$, differentiation of $-\Delta u=\lambda_Nu$
gives $-\Delta v_j=\lambda_Nv_j$. Conformal change of variables
then gives
\[
 -\Delta(v_j\circ\Phi)
 =\lambda_N|\Phi'|^2(v_j\circ\Phi)
 \quad\text{in }\mathbb D,
 \qquad \operatorname{Tr}(v_j\circ\Phi)=g_j.
\]
Uniqueness in \cref{eq:weak-potential-solution} therefore implies
\begin{equation}\label{eq:gradient-extension}
 W_g=(\nabla u)\circ\Phi,
 \qquad X_g=\nabla u,
 \qquad E(g)=0.
\end{equation}
The last equality is the equality case of the vector variational
principle.

\begin{proposition}[Multiplier identity]\label[proposition]{prop:multiplier-identity}
For every real Lipschitz function $b$ on $S^1$,
\begin{equation}\label{eq:multiplier-identity}
 E(bg)=\frac12\iint_{S^1\times S^1}
       N(s,t)(b(s)-b(t))^2g(s)\cdot g(t)\,ds\,dt.
\end{equation}
\end{proposition}
\begin{proof}
The products $bg$ and $b^2g$ are tangent Lipschitz boundary data.
Since $E$ is positive semidefinite on such data (by
\cref{prop:vector-variational}) and $E(g)=0$, for every
tangent Lipschitz $h$ and every real $t$ we have
\[
 0\le E(g+th)=2tE(g,h)+t^2E(h).
\]
If $E(g,h)\ne0$, choosing $t$ of the opposite sign and sufficiently
small contradicts this inequality. Thus $E(g,h)=0$ for all such
$h$, and in particular $E(g,b^2g)=0$. Consequently
\[
 E(bg)=E(bg,bg)-E(g,b^2g).
\]
Apply \cref{eq:boundary-energy} to both pairings on the right.
The curvature integrals cancel because their pointwise factors
are identical:
\[
 |bg|^2=g\cdot(b^2g)=b^2|g|^2.
\]

Write $b_s=b(s)$ and $g_s=g(s)$. For the harmonic integrals,
expansion gives
\[
 \begin{aligned}
 &|b_sg_s-b_tg_t|^2
    -(g_s-g_t)\cdot(b_s^2g_s-b_t^2g_t)\\
 &\qquad=(b_s^2+b_t^2-2b_sb_t)g_s\cdot g_t
          =(b_s-b_t)^2g_s\cdot g_t.
 \end{aligned}
\]
Their difference is therefore
\[
 \frac12\iint N_0(s,t)(b_s-b_t)^2g_s\cdot g_t\,ds\,dt.
\]
For the integrals with $R$, subtraction gives
\[
 \iint R(s,t)(b_t^2-b_sb_t)g_s\cdot g_t\,ds\,dt.
\]
These integrals are absolutely convergent because $b,g$ are bounded
and $R$ is integrable. Symmetry of $R$ and $g_s\cdot g_t$ permits
averaging the integrand with the one obtained by exchanging $s,t$.
Since
\[
 \frac12\bigl(b_t^2-b_sb_t+b_s^2-b_sb_t\bigr)
 =\frac12(b_s-b_t)^2,
\]
this contribution equals
\[
 \frac12\iint R(s,t)(b_s-b_t)^2g_s\cdot g_t\,ds\,dt.
\]
Adding the two contributions and using $N=N_0+R$ proves
\cref{eq:multiplier-identity}.

For completeness, the integral in that identity is absolutely
convergent. If $L=\operatorname{Lip}(b)$, then
$|b_s-b_t|\le L|s-t|$, so
\[
 \left|N_0(s,t)(b_s-b_t)^2g_s\cdot g_t\right|
 \le\frac{L^2}{\pi}\|g\|_\infty^2.
\]
This bounds the integrand off the diagonal by a constant.
For the remaining part,
\[
 \left|R(s,t)(b_s-b_t)^2g_s\cdot g_t\right|
 \le4\|b\|_\infty^2\|g\|_\infty^2R(s,t),
\]
which is integrable by \cref{eq:interaction-mass}.
\end{proof}
